% ============================================================
% CSE 8803 IUQ — Project Proposal
% Format: SIAM JUQ style (11pt, 1-inch margins, single-spaced)
% Due: March 2, 2026 at 11:59 PM EST
% ============================================================

\documentclass[11pt]{article}

% --- Page layout ---
\usepackage[margin=1in]{geometry}
\usepackage{setspace}
\singlespacing

% --- Math ---
\usepackage{amsmath, amssymb, amsthm}
\usepackage{bm}

% --- Fonts & encoding ---
\usepackage[T1]{fontenc}
\usepackage[utf8]{inputenc}
\usepackage{microtype}

% --- Figures & tables ---
\usepackage{graphicx}
\usepackage{booktabs}

% --- Hyperlinks ---
\usepackage[colorlinks=true, linkcolor=blue, citecolor=blue, urlcolor=blue]{hyperref}

% --- Bibliography ---
\usepackage{natbib}
\bibliographystyle{plainnat}

% --- Header ---
\usepackage{fancyhdr}
\pagestyle{fancy}
\fancyhf{}
\rhead{\small CSE 8803 IUQ — Project Proposal}
\lhead{\small Bo Feng}
\cfoot{\thepage}

% ============================================================

\begin{document}

% --- Title block ---
\begin{center}
  {\Large\bfseries LLM-Informed Bayesian Priors for Reducing\\
  Experimental Cost in Behavioral Economics}

  \vspace{0.8em}

  {\normalsize
    Bo Feng\\
    \texttt{bfeng66@gatech.edu}\\
    Georgia Institute of Technology\\
    CSE 8803 IUQ — Spring 2026
  }

  \vspace{0.8em}
  {\small \today}
\end{center}

\vspace{0.5em}
\hrule
\vspace{1em}

% --- Abstract ---
\begin{abstract}
Human experiments in behavioral economics are expensive, yet large
language models (LLMs) exhibit qualitatively human-like behavior in
classic economic games~\citep{horton2023homosilicus} and encode
knowledge equivalent to dozens of human observations for economic
variables~\citep{gao2026measure}.
We propose a Bayesian framework in which LLM simulations serve as an
informative prior and real human observations provide the likelihood,
keeping the two data sources in strictly separate roles to avoid
methodological circularity.
Our primary metric is the \emph{sample size reduction ratio}
$\rho = n^*_{\mathrm{LLM}} / n^*_{\mathrm{flat}}$: the minimum human
observations needed to reach a target posterior precision under the
LLM prior versus a flat prior.
We implement a flat-prior MCMC baseline and an LLM-informed SVGD
frontier method, and evaluate both on dictator-game and ultimatum-game
data from \citet{horton2023homosilicus}.
\end{abstract}

\vspace{1em}
\hrule
\vspace{1em}

% ============================================================
\section{Review of the Literature}
% ============================================================

\paragraph{\citet{horton2023homosilicus}.}
Horton introduces ``homo silicus'': LLMs queried as simulated
experimental subjects replicate qualitative findings of dictator and
ultimatum games, yielding sharing-ratio distributions broadly
consistent with published human data.
The analysis is purely qualitative with no uncertainty bounds, but it
establishes the elicitation protocol and the human-response dataset we
use as ground truth for evaluation.

\paragraph{\citet{gao2026measure}.}
Gao, Han, and Liang define the \emph{equivalent sample size} (ESS) of
an LLM's pretrained knowledge as the task-specific human observations
needed to match LLM prediction error, finding values of dozens to
hundreds for certain economic variables.
ESS is a frequentist prediction metric with no direct Bayesian
interpretation; our work operationalizes the ESS intuition inside a
Bayesian pipeline, translating it into the sample-size reduction ratio
$\rho$.

\paragraph{\citet{ludwig2025econometrics}.}
Ludwig, Mullainathan, and Rambachan show that valid LLM-assisted
inference requires strict separation of LLM training data from the
researcher's sample (``no training leakage'') and a held-out
validation sample for estimation.
Although their framework targets prediction/estimation rather than
Bayesian priors, our prior/likelihood role separation directly
implements the no-leakage requirement, and a held-out subset of the
Horton data serves as the validation sample.

\paragraph{\citet{shorinwa2025survey}.}
Shorinwa et al.\ provide a comprehensive taxonomy of UQ methods for
LLMs---covering epistemic/aleatoric decomposition, calibration
metrics, and conformal prediction---though the focus is uncertainty
\emph{of the LLM} rather than of downstream parameters $\theta$.
Their calibration framework supplies our evaluation metric (empirical
coverage of 95\% credible intervals) and their epistemic/aleatoric
distinction motivates the failure-mode sensitivity analysis.

\paragraph{\citet{gui2025challenge}.}
Gui and Toubia show that standard blind-prompting in LLM experiments
violates unconfoundedness---treatment variations alter unspecified
contextual variables, biasing simulated responses---with no general
solution when unblinding itself changes behavior.
This motivates our prior-misspecification sensitivity analysis: a
biased LLM prior proxies for blind-prompting confounding, and we
measure the resulting degradation in posterior coverage.

\paragraph{LLM-elicited priors: closest related work.}
Three concurrent works use LLMs to construct Bayesian priors; we
address each in turn.

\citet{capstick2024autoelicit} propose AutoElicit, which queries an
LLM for the mean and standard deviation of a Gaussian prior over each
feature of a \emph{linear predictive model}, then forms a
$K{=}100$-component Gaussian mixture prior and runs MCMC inference.
The evaluation metric is posterior predictive accuracy on clinical
datasets (UTI, heart disease).
Our approach differs along four dimensions.
(i)~\textit{Domain and generative model.}
AutoElicit targets regression/classification weights $\theta \in
\mathbb{R}^d$ under a linear likelihood; we target behavioral
parameters $\theta = (\alpha,\beta)$ of a Beta likelihood for sharing
ratios in economic games---a bounded, compositional parameter with no
natural Gaussian form.
(ii)~\textit{Elicitation mechanism.}
AutoElicit asks the LLM to report $(\mu_k, \sigma_k)$ directly;
we instead ask the LLM to \emph{simulate} $M{=}500$ behavioral
observations and derive $(\alpha,\beta)$ from them as
pseudo-observations, which is more natural for behavioral proportions.
(iii)~\textit{Primary metric.}
AutoElicit measures predictive accuracy at fixed $n$; our primary
metric is the \emph{sample size reduction ratio}
$\rho = n^*_{\mathrm{LLM}}/n^*_{\mathrm{flat}}$---a quantity not
reported by any prior work.
(iv)~\textit{ESS-calibrated prior weight.}
AutoElicit assigns no principled weight to the LLM prior; we scale
pseudo-observation count by $w$ calibrated to the equivalent sample
size of \citet{gao2026measure}, grounding prior strength in a
measurable quantity.
AutoElicit also does not adopt SVGD, the no-leakage framework of
\citet{ludwig2025econometrics}, or a systematic prior-misspecification
analysis.

\citet{gouk2024autoprior} (described in \citealt{capstick2024autoelicit})
use LLMs to generate synthetic feature--label pairs from public
datasets, which implicitly define a prior through a separate likelihood
term during training.
This is \emph{not} a Bayesian prior in the usual sense: LLM-generated
data enters the likelihood, not the prior, so there is no
prior/likelihood role separation.
\citeauthor{capstick2024autoelicit} note this makes the method
susceptible to LLM data memorization on public benchmarks; our
held-out Horton data and explicit role separation avoid this failure
mode.

\citet{huang2025llmbi} propose LLM-BI, a proof-of-concept pipeline in
which an LLM translates a \emph{user's natural-language belief
description} into prior distributions and a likelihood function for
Bayesian linear regression.
LLM-BI validates that posteriors from LLM-specified priors agree with
posteriors from manually specified priors on synthetic data.
Our work differs on three points.
(i)~LLM-BI relies on the \emph{user} to supply prior beliefs in
natural language; we rely on the \emph{LLM's own behavioral knowledge}
from pretraining, without any user-provided prior information.
(ii)~LLM-BI is evaluated on synthetic data with no human experimental
data; our evaluation uses real dictator- and ultimatum-game data.
(iii)~LLM-BI does not address sample-size reduction, prior
calibration, training-data leakage, or failure-mode robustness.

% ============================================================
\section{Description of the Proposed Project}
% ============================================================

\subsection{Problem Statement and Significance}

We consider Bayesian inference over human behavioral parameters in economic
experiments, with an LLM serving as an informative prior.
Let $\theta \in \Theta$ parametrize the distribution of human behavior---in
the dictator game, $\theta = (\alpha, \beta)$ are the shape parameters of a
Beta likelihood $p(y \mid \theta) = \mathrm{Beta}(y; \alpha, \beta)$ for
sharing ratios $y \in [0,1]$.
Given $n$ human observations $\mathbf{y} = \{y_1, \ldots, y_n\}$, the
posterior is:
%
\begin{equation}
  p(\theta \mid \mathbf{y}) \;\propto\; p(\mathbf{y} \mid \theta)\cdot p(\theta).
  \label{eq:bayes}
\end{equation}
%
Our quantity of interest (QoI) is the posterior $p(\theta \mid \mathbf{y})$
and, operationally, the minimum sample size $n^*$ required to reach a target
posterior precision $\sigma^2_{\mathrm{target}}$.
The central UQ question is: \textit{how much does an LLM-informed prior
reduce $n^*$ relative to a flat prior?}

Human experiments are costly (\$10--50 per participant on MTurk), yet
\citet{gao2026measure} show LLM knowledge is worth dozens to hundreds of
human observations for certain economic variables.
Our framework assigns LLM predictions and human data to \emph{strictly
separate} roles in \eqref{eq:bayes}: LLM outputs enter only through the
prior, human data only through the likelihood.
This avoids the circularity of calibrating an LLM against human data and
then using it to predict human behavior---the role-separation also
directly implements the no-leakage requirement of \citet{ludwig2025econometrics}.
One caveat: a modern LLM trained after January 2023 may have seen
Horton's published results; we treat this as a failure mode and address it
in the sensitivity analysis.

\subsection{Methodology}

\paragraph{Prior construction.}
For each experimental condition $x$, we query GPT-4o
(\texttt{gpt-4o-2024-08-06}) with a natural-language description of the game
and collect $M = 500$ simulated sharing ratios $\{\tilde{y}_1, \ldots,
\tilde{y}_M\}$.
We treat these as $M$ pseudo-observations and use them as sufficient
statistics to initialize a conjugate Beta prior over $\theta = (\alpha,
\beta)$, scaled by an effective weight $w \in (0,1]$ calibrated to the
equivalent sample size of \citet{gao2026measure}.
This yields the LLM-informed prior $p_{\mathrm{LLM}}(\theta \mid x)$.

\paragraph{Baseline.}
A flat prior $p(\theta) \propto 1$ combined with MCMC (Metropolis--Hastings)
provides the posterior $p(\theta \mid \mathbf{y})$.
We estimate the minimum sample size $n^*_{\mathrm{flat}}$ by running
inference at $n \in \{5, 10, 20, 50, 100, 200\}$ and interpolating the
threshold where posterior variance on $\mu = \alpha/(\alpha+\beta)$ drops
below $\sigma^2_{\mathrm{target}} = 0.1 \times \mathrm{Var}_{\mathrm{prior}}(\mu)$.

\paragraph{Frontier method.}
We substitute $p_{\mathrm{LLM}}(\theta)$ for the flat prior and use Stein
Variational Gradient Descent (SVGD) for posterior inference.
Although the present problem is low-dimensional ($\theta \in \mathbb{R}^2$),
SVGD is adopted as the frontier method because it generalizes naturally to
multiple experimental conditions and higher-dimensional behavioral parameter
spaces; the particle-based representation also makes posterior geometry
immediately visualizable.
We obtain $n^*_{\mathrm{LLM}}$ by the same interpolation procedure.

\paragraph{Evaluation.}
Our primary metric is the \textbf{sample size reduction ratio}
$\rho = n^*_{\mathrm{LLM}} / n^*_{\mathrm{flat}}$,
where $\rho < 1$ indicates cost reduction.
We additionally evaluate:
\begin{itemize}
  \item \textit{Posterior calibration}: empirical coverage of 95\%
        credible intervals on a held-out split of the Horton data.
  \item \textit{Compute--accuracy tradeoff}: wall-clock time and ESS
        for SVGD vs.\ MCMC at matched accuracy.
  \item \textit{Failure mode}: we artificially shift the LLM prior mean
        by $\pm 1, 2$ standard deviations and measure coverage degradation,
        identifying conditions under which the informative prior is harmful.
\end{itemize}
The framework is applied to the dictator game and ultimatum game from
\citet{horton2023homosilicus}.

\paragraph{Compute plan and risks.}
LLM queries require $M \times K \approx 500 \times 4$ API calls (two games,
two conditions each), costing under \$5.
MCMC and SVGD runs are CPU-only and complete in minutes for a 2-parameter
model.
The primary risk is LLM prior misspecification (e.g., due to training-data
leakage from Horton's published results); this is explicitly studied in the
failure-mode evaluation above.

\subsection{Contributions}

We contribute a Bayesian framework that incorporates LLM simulations as
informative priors with strict prior/likelihood separation, an empirical
characterization of sample-size reduction $\rho$ across multiple game
settings, and a robustness analysis identifying when LLM priors help
versus hurt posterior inference.

\subsection{Plan and Timeline}

\begin{table}[h]
  \centering
  \small
  \begin{tabular}{llp{7cm}}
    \toprule
    \textbf{Date} & \textbf{Milestone} & \textbf{Deliverable} \\
    \midrule
    Mar 2   & Proposal submission  & This document + 3--4 slides \\
    Mar 11  & Check-in 1           & Baseline implemented (flat prior + MCMC) on
                                     Horton (2023) dictator game; initial
                                     posterior results \\
    Apr 8   & Check-in 2           & LLM prior construction + SVGD implemented;
                                     preliminary comparison to baseline \\
    Apr 27  & Final presentation   & 15-min talk with full evaluation results \\
    Finals  & Final report         & 10-page report + reproducible code \\
    \bottomrule
  \end{tabular}
  \caption{Project timeline.}
  \label{tab:timeline}
\end{table}

% ============================================================
% References
% ============================================================

\bibliography{references}

\end{document}
